تفاوت دو رویکرد بر اساس معیارهای مطرح شده:

\begin{enumerate}
    \item {\lr{Coupling}:}
    \begin{itemize}
        \item {همگام:} دارای \lr{Coupling} بالایی است. سرویس سفارش باید مستقیما آدرس (\lr{IP/URL}) و ساختار سرویس‌های پرداخت و ایمیل را بشناسد.
        \item {ناهمگام:} دارای \lr{Loose Coupling} است. سرویس سفارش فقط رویداد را به \lr{Message Broker} می‌فرستد و هیچ شناختی از گیرنده‌ها ندارد.
    \end{itemize}

    \item {\lr{Latency}:}
    \begin{itemize}
        \item {همگام:} تاخیر کل برابر است با مجموع زمان پاسخ‌دهی تمام سرویس‌ها. سرویس سفارش باید منتظر بماند تا پرداخت و ایمیل کارشان تمام شود که سرعت را پایین می‌آورد.
        \item {ناهمگام:} تاخیر برای کاربر به شدت کاهش می‌یابد؛ زیرا به محض ثبت سفارش و ارسال ایونت، پاسخ به کاربر برگشت داده می‌شود و بقیه کارها در پس‌زمینه انجام می‌شود.
    \end{itemize}

    \item {\lr{Fault Tolerance}:}
    \begin{itemize}
        \item {همگام:} پایین است. اگر سرویس ایمیل یا پرداخت در آن لحظه در دسترس نباشد، کل فرآیند ثبت سفارش با خطا مواجه می‌شود.
        \item {ناهمگام:} بالا است. اگر سرویس ایمیل خراب باشد، پیام در \lr{Queue} یا \lr{Broker} باقی می‌ماند و به محض بالا آمدن سرویس، پردازش می‌شود؛ بدون اینکه تراکنش اصلی خراب شود.
    \end{itemize}

    \item {\lr{Scalability}:}
    \begin{itemize}
        \item {همگام:} به سختی مقیاس می‌شود؛ زیرا فشار کاری مستقیما به سرویس‌های بعدی منتقل می‌شود و در صورت هجوم کاربران، گلوگاه ایجاد می‌شود.
        \item {ناهمگام:} بسیار بالا است. \lr{Message Broker} مانند یک ضربه‌گیر عمل کرده و پیام‌ها را ذخیره می‌کند؛ سرویس‌ها می‌توانند متناسب با توان خود و به صورت ناهمزمان پیام‌ها را مصرف کنند.
    \end{itemize}
\end{enumerate}

با توجه به تحلیل بالا، بهترین استراتژی ترکیب هر دو رویکرد به شکل زیر است:

\begin{enumerate}
    \item {ارتباط با سرویس پرداخت با رویکرد همگام:}

فرآیند پرداخت یک عملیات حیاتی و وابسته به زمان است. سرویس سفارش باید فورا متوجه شود که آیا پرداخت با موفقیت انجام شد یا خیر تا وضعیت سفارش را مشخص کند. بنابراین استفاده از یک فراخوانی همگام مانند \lr{REST} در اینجا منطقی‌تر است.

    
    \item {ارتباط با سرویس ارسال ایمیل با رویکرد ناهمگام:}

    ارسال ایمیل تاییدیه خرید، تأثیری روی موفقیت آمیز بودن خود فرآیند ثبت سفارش ندارد. کاربر نباید معطل ارسال ایمیل بماند. بنابراین، سرویس سفارش پس از نهایی شدن پرداخت، رویدادی مانند \lr{OrderConfirmed} را در \lr{Message Broker} منتشر می‌کند و سرویس ایمیل آن را به صورت ناهمگام برمی‌دارد و پردازش می‌کند.
\end{enumerate}

\end{document}